\section{Is Emotion Legible From a Face at All?}
\label{sec:11.7}
Whether a face reveals what a person feels is not settled. Facial-expression-recognition AI is built almost entirely on Paul Ekman and Wallace Friesen's claim that a small set of emotions — anger, disgust, fear, happiness, sadness, surprise — is universally expressed and universally readable, based on the Fore study: roughly 85 percent cross-cultural agreement on which face matched which emotion story, among people with almost no prior exposure to Western media \autocite{ekman1971constants}. What that evidence settles about what emotion is has already been asked, and the dispute is open. The question here is the narrower one of whether a face is readable.

A meta-analysis spanning both face and voice found real above-chance recognition across cultures, but also a consistent “in-group advantage”: people read emotion more accurately from members of their own national, ethnic, or linguistic group than from anyone else's, an asymmetry that shrinks, but does not vanish, the more contact the groups have with each other \autocite{elfenbein2002universality}. A system trained on one population's expressive norms inherits that same asymmetry, and will misread everyone outside it.

In 2019 the underlying claim collapsed further. The Association for Psychological Science commissioned five researchers from competing theoretical camps to review the evidence; the assignment was expected to take a few months and took two years, working through more than a thousand studies before the panel reached a unanimous conclusion: a person's emotional state cannot be reliably inferred from facial movement alone, in the strong sense the six-emotion model requires \autocite{barrett2019emotional}. That verdict is not a complaint about noisy classifiers needing more training data. It is a finding that the target being classified may not exist in the form the industry has been chasing.

This result has already reshaped a deployed product, which section~\ref{sec:12.1} works through. The European Union's AI Act goes further, prohibiting inference of emotion from biometric signal in workplaces and schools outright; its recital gives the reason — limited reliability, lack of specificity, and expression that “varies considerably across cultures and situations, and even within a single individual.” One of the Act's few bans rests on the science failing, not on how the technology might be misused \autocite{eu2024aiact}.

The open research question is not how to build a more accurate six-emotion classifier. Nor is it quite the skeptical form of the same question — whether any face-to-feeling mapping survives once individual and cross-cultural variation are taken seriously, under what constraints, and against what error rate. Both of those ask how good the instrument is, and section~\ref{sec:2.2.1} gives the reason to ask something else. A face-reading pipeline registers every face against a reference face before it computes anything, and the reference was chosen by somebody; what comes out the far end as an error rate is a distance from it. So the question is what the instrument is measuring distance from, who put that there, and what a deployment does with the ordering it gets back. The instrument is not a neutral way of finding out whether emotion is legible from a face. It is the politics, already built.

\runin{The near-term work}
The operational form of that is an audit of the registration step rather than of the output, and it can be put to any system already built. Publish what the system aligns to — the reference model, or the summary statistics of the corpus it was fit on — and report error disaggregated by distance from that reference, rather than only by the named group of the person being read. Two things follow that a group-wise table does not supply. Distance is continuous, so the claim becomes falsifiable in a way a comparison between categories is not: if error does not climb with distance from the reference, the account in section~\ref{sec:2.2.1} is wrong about this system. And a published reference is contestable in a way an error rate is not. A vendor can accept a disparity finding and answer it with more data; a vendor cannot answer a published reference face except by choosing a different one and saying which.

The older measurement keeps its place, because it is the part of the accuracy question already shaped like a relation rather than a property. Report a classifier's error against the population whose expressive norms it was trained on alongside its error against everyone else, and publish the in-group advantage as a number belonging to that system. The falsifier there runs in the system's favor for once: flat error rates across groups would mean it has done something the meta-analytic evidence says is unavailable. The data is the part nobody has built, in both cases. An in-group advantage is a relation between reader and read, and the standard corpora record only the subject; a reference model is a thing vendors have never had a reason to publish. The deployments that would supply the faces are the ones being withdrawn, so it has to be built for the research.
